% Generated draft: requires mathematical and source review.
\documentclass[11pt]{article}
\usepackage[margin=25mm]{geometry}
\usepackage{fontspec}
\setmainfont{Latin Modern Roman}
\usepackage{amsmath,amssymb,mathtools,unicode-math}
\setmathfont{Latin Modern Math}
\usepackage{enumitem,needspace,xurl,hyperref}
\setlength{\parindent}{0pt}
\setlength{\parskip}{5pt}
\setlength{\emergencystretch}{3em}
\newcommand{\N}{\mathbb{N}}
\newcommand{\Z}{\mathbb{Z}}
\newcommand{\Q}{\mathbb{Q}}
\newcommand{\R}{\mathbb{R}}
\newcommand{\C}{\mathbb{C}}
\newcommand{\F}{\mathbb{F}}
\newcommand{\A}{\mathbb{A}}
\newcommand{\PP}{\mathbb{P}}
\newcommand{\E}{\mathbb{E}}
\newcommand{\eps}{\varepsilon}
\newcommand{\dd}{\,\mathrm{d}}
\begin{document}
\section*{UnsolvedMath ID: 10600001}
\textbf{Source problem code:} catalogue entry
\subsection{Definitions and mathematical statement}
A closed virtual knot is represented by a circle diagram with classical over/under crossings and virtual crossings, modulo the virtual Reidemeister and detour moves. Let \(\mathcal V\) be its set of equivalence classes and \(U\in\mathcal V\) the unknot. Let \(\mathcal L\) be the set of long virtual knot classes: diagrams on an oriented line, standard outside a compact set, with equivalence preserving the ends. Long concatenation \(L_1*L_2\) is ordered and well-defined, and closure gives \(\operatorname{cl}:\mathcal L\to\mathcal V\). The source stresses that connected sum of closed virtual knots is not a well-defined operation on \(\mathcal V\): choices of diagrams and connecting sites matter.

\emph{Compound problem: recognizing the Kishino knot and classifying the connected-sum phenomenon.} The source records the Kishino virtual knot as nontrivial despite a diagrammatic decomposition into two trivial virtual knots. Other proofs of its nontriviality may be illuminating, but the supplied record does not present nontriviality itself as an unsolved detection question. The remaining classification request is: classify when choices of representatives and connecting sites for two closed virtual unknots give a nontrivial connected sum. If \(D_1,D_2\) represent \(U\) and \(a\) records the admissible connecting choices, determine when
\[
 [D_1\mathbin{\#_a}D_2]\neq U,
\]
and classify the resulting phenomenon in the intended equivalence sense. The source does not specify whether the requested classification is of outcomes, decomposition data, or a necessary and sufficient criterion on that data.

The source calls this closely related to classifying the long virtual knots whose closures are the unknot, namely
\[
 \operatorname{cl}^{-1}(U)=\{L\in\mathcal L:\operatorname{cl}(L)=U\}.
\]
In particular, the long-knot framework can record ordered pairs with \(\operatorname{cl}(L_1)=\operatorname{cl}(L_2)=U\) but \(\operatorname{cl}(L_1*L_2)\neq U\). This expresses the relevant choice-dependent mechanism; the source does not assert an equivalence of complete classification problems or that such \(L_i\) are trivial as long knots.

The supplied discussion lists detection by the Jones polynomial of the \(3\)-strand cable, the \(\Xi\)-polynomial, the quaternionic biquandle, and the surface bracket polynomial of Dye and Kauffman. It credits Kadokami with a proof using the immersion class of a shadow curve in genus two. It also says that Manturov's parity bracket makes the flat Kishino diagram essentially its own invariant and exhibits nontriviality. The exact diagram, cabling convention, invariant definitions, and several intended citations are absent.
\subsection{Short English statement}
The Kishino knot is already recognized as nontrivial by several supplied methods. The classification problem asks how two closed virtual unknots can give a nontrivial connected sum when connecting choices matter, and how this relates to nontrivial long virtual knots with trivial closures.
\subsection{Sources}
\begin{itemize}
\item Local UnsolvedMath registry, permanent ID 10600001, source code catalogue entry, ``Virtual-knot problem 1 --- Recognising the Kishino Knot.'' No individual proposer is supplied.
\item \emph{Unsolved Problems in Virtual Knot Theory and Combinatorial Knot Theory}, revised survey, problem-list item 1. Supplied URL: \url{https://arxiv.org/abs/1409.2823}; recorded access 2026-07-29, extraction source-tex. The dedicated section has 60 top-level enumerated items, with compound items retained intact.
\item The supplied text names Dye and Kauffman for the surface bracket polynomial, Kadokami for the genus-two shadow-curve proof, and Manturov for the parity-bracket approach. References for several listed detection methods and the exposition of the parity proof are missing from the extracted text; no bibliographic details are invented.
\item Supplied status partially\_solved and PARTIAL-PROGRESS. The summary marks nontriviality SOLVED, with multiple proofs before 2014 and later, and the connected-sum classification PARTIAL, saying no complete classification was found and the specific question appears open despite decomposition literature. Background requires a current primary-literature status audit and review of redistribution terms for the publicly accessible author/publisher copy, both NEEDS\_REVIEW.
\end{itemize}
\end{document}
